% Related Work chapter — Parkinson's detection from radar micro-Doppler gait.
% Requires references.bib (cite keys used below). Compiles with \bibliographystyle
% {IEEEtran} or a natbib style; \cite is used throughout for portability.
% Convert to Word if needed: pandoc RELATED_WORK.tex --bibliography=references.bib -o related_work.docx

\chapter{Related Work}
\label{ch:related}

This chapter reviews the literature at the intersection of contactless radar
sensing, micro-Doppler signal processing, and machine-learning gait analysis,
and positions the present work within it. The survey is organised from the
physical substrate upward: the case for radar as a gait sensor
(Section~\ref{sec:contactless}); what a radar micro-Doppler signature physically
encodes and how reliably it can be measured
(Section~\ref{sec:foundations}); how deep learning has been applied to
micro-Doppler spectrograms (Section~\ref{sec:dl}); strategies for learning from
the small datasets typical of this field (Section~\ref{sec:small}); the nascent
use of radar for neurodegenerative disease
(Section~\ref{sec:neuro}); and the validation pitfalls that recur throughout
(Section~\ref{sec:pitfalls}). Section~\ref{sec:gap} distils the resulting
research gap that this thesis addresses.

\section{Contactless gait sensing}
\label{sec:contactless}

Quantitative gait analysis is clinically valuable for detecting and monitoring
movement disorders, but the dominant instrumented modalities each carry a cost.
Wearable inertial sensors are accurate and inexpensive yet must be worn,
constraining continuous use and biasing natural gait; optical motion capture and
video are accurate but expensive, illumination-dependent, and privacy-invasive.
Radar has emerged as a complementary alternative that is contactless,
insensitive to lighting and clothing, and privacy-preserving, making it
attractive for continuous and in-home monitoring
\cite{seifert2020,gurbuz2024,alanazi2022}. These properties motivate the growing
body of work that uses radar to characterise human gait, which the remainder of
this chapter reviews.

\section{Radar micro-Doppler foundations and gait-parameter validation}
\label{sec:foundations}

When a radar illuminates a walking person, different body parts moving at
different radial velocities induce distinct Doppler shifts; the resulting
time--frequency representation, the micro-Doppler signature, encodes the slow,
steady motion of the torso at low Doppler frequencies and the faster swings of
the limbs at higher frequencies \cite{seifert2020}. Establishing what of this
signature is physically trustworthy has been the subject of careful validation
studies. Seifert et al.\ \cite{seifert2020} benchmarked a 24\,GHz Doppler radar
against marker-based motion capture and instrumented-treadmill ground-reaction
forces, showing that radar recovers spatiotemporal rhythm and pace parameters
(stride and step time, cadence, stride length) reliably, but that joint-velocity
estimates and gait-event timing are strongly dependent on geometry, viewing
angle, and walking speed. Complementing this, the survey of Gurbuz et al.\
\cite{gurbuz2024} catalogues radar-based gait-parameter estimation for fall-risk
assessment and notes a crucial limitation: radar captures cadence and step-time
well but measures \emph{gait variability}---the stride-to-stride fluctuation
that is a hallmark of parkinsonian gait---only poorly to moderately. The radar
network that produced the dataset used in this thesis, López-Delgado et al.\
\cite{lopezdelgado2026}, sits in this validation tradition: it compares
multi-node radar configurations against a Vicon reference and demonstrates
accurate spatiotemporal parameter extraction in both healthy and Parkinson's
populations, reporting that foot-oriented nodes best recover temporal parameters
in motorically impaired subjects and that the torso-velocity--heel-strike
coupling reliable in healthy gait \emph{breaks down} in Parkinson's disease.
Critically, however, that work stops at parameter measurement and contains no
disease classifier---the point of departure for the present thesis.

\section{Deep learning for micro-Doppler gait classification}
\label{sec:dl}

A second line of work treats the micro-Doppler spectrogram as an image or a
sequence and applies machine learning to classify gait. Two architectural
families dominate. The first feeds the full time--frequency spectrogram to a
convolutional neural network (CNN); the second extracts one-dimensional velocity
time-series (for example the upper, lower, and power-weighted Doppler envelopes)
and models them with recurrent networks. Hayashi/Saho et al.\
\cite{hayashi2021} exemplify the latter, classifying young versus elderly gait
from 24\,GHz micro-Doppler with a long short-term memory network on extracted
envelopes, reaching 94.9\% accuracy and reporting that the stance-phase
(lower-envelope) feature alone was most discriminative. A companion simulation
study \cite{hoshiga2021} corroborates that leg-motion features carry most of the
age-discriminative signal. CNN-based approaches on spectrograms are equally
common across human-activity and gait recognition
\cite{papanastasiou2020,ni2020,fard2026,nguyen2024,ha2023}, with reported
accuracies frequently in the 90--98\% range. Notably, Papanastasiou et al.\
\cite{papanastasiou2020} and Ni and Huang \cite{ni2020} show that micro-Doppler
encodes \emph{individual identity} strongly enough to recognise people
(>93.5\% and 96.7\% respectively)---a result that, as Section~\ref{sec:pitfalls}
argues, has direct and cautionary implications for how disease classifiers must
be validated. Alternative front-ends have also been explored, including treating
the Doppler signature with audio-processing pipelines \cite{ha2023} and
low-cost millimetre-wave sensing for gait \cite{alanazi2022}.

\section{Learning from small radar datasets}
\label{sec:small}

Radar gait datasets are typically small---tens of subjects---because data
collection requires supervised sessions with specialised hardware. Two
strategies recur for coping with data scarcity. The first is transfer learning:
Seyfio\u{g}lu et al.\ \cite{seyfioglu2018} show that deep networks pre-trained on
diversified micro-Doppler transfer effectively to new motion-classification
tasks, and Park et al.\ \cite{park2016} demonstrated transfer learning on
micro-Doppler as early as 2016. The second is synthetic data augmentation:
Erol et al.\ \cite{erol2020} use a kinematically sifted auxiliary-classifier
GAN to synthesise realistic radar micro-Doppler signatures and expand small
training sets. Both strategies are directly relevant to a 58-subject study and
inform the transfer-learning and augmentation choices adopted in this thesis.

\section{Radar for neurodegenerative disease and Parkinson's}
\label{sec:neuro}

Applying radar gait analysis specifically to neurodegenerative disease is
recent and sparse. Cai et al.'s STRIDE system \cite{cai2023} uses radar
intelligence for Alzheimer's-disease-and-related-dementia risk evaluation,
illustrating radar's promise for passive cognitive-decline screening. For
Parkinson's disease specifically, machine-learning gait classification has been
pursued mainly through non-radar modalities such as force sensors and wearables
\cite{neurosurgery2025}, and the radar literature on Parkinson's remains largely
confined to parameter validation \cite{lopezdelgado2026} rather than automated
diagnosis. To the best of our knowledge, no published work performs
subject-independent Parkinson's-versus-control classification directly from
radar micro-Doppler spectrograms---the specific gap this thesis targets.

\section{Validation pitfalls: subject leakage, artifacts, and confounds}
\label{sec:pitfalls}

The high accuracies reported across the micro-Doppler classification literature
must be read against a set of recurring methodological hazards, because several
are known to inflate results. First, because micro-Doppler encodes individual
identity so strongly \cite{papanastasiou2020,ni2020}, any evaluation that allows
recordings from the same subject to appear in both training and test sets
permits the model to recognise \emph{people} rather than the target condition;
subject-independent (leave-one-subject-out) validation is therefore essential,
yet much of the field relies on random hold-out splits
\cite{hayashi2021,hoshiga2021}. Second, spectrogram CNNs are prone to keying on
nuisance artifacts rather than gait: Hayashi/Saho et al.\ \cite{hayashi2021}
explicitly demonstrate that an AlexNet reaching 97.8\% did so by exploiting
background-noise differences arising because their two groups were recorded at
different locations, not by learning gait---a textbook case of shortcut
learning. Third, group-level confounds are pervasive. Age is especially
insidious for a Parkinson's study: because micro-Doppler alone separates young
from elderly walkers at 94.9\% \cite{hayashi2021} and Parkinson's cohorts are
typically older than controls, a classifier can appear to detect disease while
in fact detecting age. Finally, the very features most diagnostic of
parkinsonian gait---reduced velocity (bradykinesia) and increased stride-to-stride
variability---are among those radar measures least reliably \cite{gurbuz2024},
so pipelines that aggressively normalise amplitude or use short temporal windows
risk discarding the pathology they seek. Together these observations argue that
methodological rigour---subject-independent validation, explicit confound
control, and interpretability checks confirming the model attends to
physiologically plausible Doppler bands---is the central unsolved problem in this
niche, not raw accuracy.

\section{Research gap and positioning of this thesis}
\label{sec:gap}

The reviewed literature supports three conclusions. First, radar micro-Doppler
is a validated, information-rich substrate for gait analysis, but the components
of the signal that are trustworthy (rhythm, pace) differ from those most
diagnostic of Parkinson's (variability, absolute velocity)
\cite{seifert2020,gurbuz2024}. Second, deep learning on micro-Doppler is mature
for coarse tasks such as activity recognition, person identification, and
young-versus-elderly discrimination, but has not been applied to
subject-independent Parkinson's classification
\cite{hayashi2021,papanastasiou2020,ni2020,lopezdelgado2026}. Third, the
field's headline accuracies are repeatedly undermined by subject leakage,
shortcut learning on recording artifacts, and uncontrolled confounds---above all
age \cite{hayashi2021}.

This thesis addresses precisely this gap. It performs, to our knowledge, the
first subject-independent deep-learning classification of Parkinson's disease
versus healthy controls from radar micro-Doppler gait spectrograms, using the
multi-node FMCW dataset of \cite{lopezdelgado2026}. Methodologically, it adopts
the safeguards the literature identifies as necessary but rarely applies
together: leave-one-subject-out cross-validation with subject-level score
aggregation to eliminate identity leakage; fixed-length windowing to remove a
trial-duration confound by construction; explicit duration-matched and
age-matched control analyses to guard against the age confound highlighted by
\cite{hayashi2021}; label-permutation testing and bootstrap confidence intervals
for honest uncertainty reporting; and gradient-based attention checks to verify
that the network attends to the foot and torso Doppler bands rather than
artifacts. In doing so, the work aims to establish not merely whether Parkinson's
gait is detectable from radar, but whether it is detectable \emph{under
validation conditions strong enough to be trusted}.
